(1) Bloch变换 $|k\rangle\propto\sum_ne^{ikna}|n\rangle$ 使能量 $E(k)=-2J\cos ka$，$k\in[-\pi/a,\pi/a]$。$v_g=\hbar^{-1}\partial_kE=2Ja\sin ka/\hbar$，最大速度为 $2Ja/\hbar$。

(2) 位置传播核是 $\psi_n(t)=\int_{-\pi}^{\pi}(d\theta/2\pi)e^{in\theta}e^{i(2Jt/\hbar)\cos\theta}=i^nJ_n(2Jt/\hbar)$，在 $t=0$ 有 $\psi_n(0)=\delta_{n0}$。于是概率 $p_n=J_n(z)^2$、$z=2Jt/\hbar$。

(3) 由离散Fourier幺正性 $\sum_nJ_n(z)^2=1$，故总概率始终为1。由于 $J_{-n}(z)^2=J_n(z)^2$，$\langle n\rangle=0$。Bessel矩恒等式 $\sum_nn^2J_n(z)^2=z^2/2$ 给 $\boxed{\langle x^2\rangle=a^2z^2/2=2a^2(Jt/\hbar)^2}$。

(4) $z\ll1$ 时 $J_0\simeq1-z^2/4$、$J_{\pm1}\simeq\pm z/2$，所以 $p_{\pm1}\simeq (Jt/\hbar)^2$，中心概率 $\simeq1-2(Jt/\hbar)^2$。长时分布显著传播宽度 $\sqrt{\langle x^2\rangle}\propto t$，属于弹道传输；经典连续时间无偏Markov随机行走则 $\langle x^2\rangle\propto t$，因为量子相干路径幅度相加而非概率直接相加。
